\subsection{二分与完整枚举}

记 $J=|A|$。实现先计算一次样本参考，再在降序网格中寻找最大的可接纳下标：
\par\noindent\begin{minipage}{\linewidth}
\begin{enumerate}[leftmargin=*,itemsep=0.15em]
\item 置 $l=0$、$h=J-1$。下标零为参考回退。
\item 当 $l<h$，测试 $j=\lfloor(l+h+1)/2\rfloor$，运行 $A[j]$ 并计算 $D_S(A[j])$。
\item 若 $D_S(A[j])\leq\epsilon$，置 $l=j$；否则置 $h=j-1$。
\item 返回 $A[l]$；调用方在处理复用批次前应用该值。
\end{enumerate}
\end{minipage}\par\medskip
八点网格除参考外至多需要三次完整候选预算搜索。调用语义、主机比较复杂度和计时边界见\cref{sec:a-budget-contract}。

完整枚举重新计算参考，在完全相同样本上测试其他七个预算，选择最小被接纳乘数。两种计时都含新参考，统一成本边界。即使接纳集合存在空洞，枚举仍评价整个有限网格；二分要保证找到相同最小值，则需要接纳下标构成前缀。

\subsection{单调性的充分条件}

\par\noindent\begin{minipage}{\linewidth}
\begin{proposition}[嵌套候选保留参考成员]
\label{prop:6-monotonicity}
对单条查询，假定路由固定，且 $\alpha\leq\beta\leq\alpha_{\rm ref}$ 时
$C_\alpha\subseteq C_\beta\subseteq C_{\alpha_{\rm ref}}$。假定至少有 $k$ 个有效候选、精排分数固定、最终并列使用同一个确定全序。每个结果均为该顺序下精确 top-$k$，则
\begin{equation}
 |\mathcal R_\alpha\cap\mathcal R_{\rm ref}|
 =|C_\alpha\cap\mathcal R_{\rm ref}|,
 \qquad D_S(\alpha)\geq D_S(\beta).
 \label{eq:6-monotonicity}
\end{equation}
后一结论要求条件对每个采样查询均成立。
\end{proposition}
\begin{proof}
对$x\in\mathcal R_{\rm ref}\cap C_\alpha$，共同全序给出
\begin{equation*}
 |\{y\in C_\alpha:y\prec x\}|
 \leq|\{y\in C_{\alpha_{\rm ref}}:y\prec x\}|<k.
\end{equation*}
故$x\in\mathcal R_\alpha$。结合$\mathcal R_\alpha\subseteq C_\alpha$，
即得\cref{eq:6-monotonicity}的交集恒等式。嵌套性给出
$|C_\alpha\cap\mathcal R_{\rm ref}|
 \leq|C_\beta\cap\mathcal R_{\rm ref}|$。
对查询求和缺失参考数，即得$D_S(\alpha)\geq D_S(\beta)$。
\end{proof}
\end{minipage}\par\medskip

在同一固定完整 key 下进行精确主码 top-$c$，可在 $c$ 增大时提供嵌套前提。再固定精排完整 key，接纳便单调，下标形成前缀，二分找到网格中最小可接纳预算。GT recall 和执行速度均不在定理中：近似精排分数可能随候选池扩大而提升错误近邻的名次，运行时间也依赖硬件。

\begin{figure}[tbp]
\centering
% Synthetic sets only: all primary scores can tie; refined order is a<b<c.
\begingroup
\small\setlength{\tabcolsep}{7pt}
\begin{tabular}{@{}lcccc@{}}
\toprule
Case & $C_1$ & $C_2$ & $C_3$ & Reference overlap \\
\midrule
Nested & $\{b\}$ & $\{b,c\}$ & $\{a,b,c\}$ & $0,\ 0,\ 1$ \\
Non-nested & $\{a\}$ & $\{b,c\}$ & $\{a,b,c\}$ & $1,\ 0,\ 1$ \\
\bottomrule
\end{tabular}
\endgroup

\caption[嵌套与非嵌套候选池]{构造的 $k=1$ 例子，精排顺序 $a\prec b\prec c$，参考为 $\{a\}$。嵌套情形保持参考成员；主码分数全并列时，不同并列选择可能产生非嵌套的中间候选池，破坏单调性。这些是解释用集合，不是 GPU 实测输出。}
\label{fig:6-nesting}
\end{figure}

\subsection{评测分组路径的推论}

运行头源码默认精确主码选择。同一次校准抽取内，$S$、路由和扫描线程数固定，\texttt{scheduled\_groups(S)} 与 $c$ 无关。改变预算影响缓冲容量和接纳工作，分组评分器仍使用相同 FP32 LUT 项与累加顺序。精排固定 $B_r$ 和 16 lanes，保持每候选的残差归约；最终选择使用 $(\refineddist,\mathrm{ID})$。

\par\noindent\begin{minipage}{\linewidth}
\begin{corollary}[分组校准]
\label[corollary]{cor:6-grouped-calibration}
对评测中的分组配置，在第四、五章的表示与共享界契约下，各预算的主码保留集合嵌套。因此\Cref{prop:6-monotonicity}适用，二分与枚举返回同一最小被接纳网格预算。
\end{corollary}
\begin{proof}
\Cref{lem:5-buffer,prop:5-local-coverage,prop:5-shared-bound}给出共同完整键 $(\primarydist,\mathrm{position})$ 下的精确主码 top-$c$，其前缀随 $c$ 增大而嵌套。固定的逐候选精排与精确最终全序提供\Cref{prop:6-monotonicity}的其余前提。
\end{proof}
\end{minipage}\par\medskip

推论覆盖$S=32,64,128$的2,304次抽取中的2,176次；其余128次是
OpenAI-3072在$S=128$的距离单block路径。补证逐次记录完整网格计数，
2,304条曲线全部单调，决策均与枚举一致（\cref{sec:7-budget-trace}）。

\paragraph{全池覆盖范围。}
十二次$S=500$调用另计：两个Vogue单元分组，十个使用$\gscan=0$。
补证新增全池枚举，十二条曲线均单调，所选预算均等于枚举得到的最小接纳预算，
因此关闭了有限网格的实测缺口。十个距离路径仍不在
\Cref{cor:6-grouped-calibration}中；要给一般嵌套保证，仍可要求共同截断顺序，
或所有网格截断（含参考）均无跨界并列。匹配留出路径有助迁移，但不证明一般条件。

\paragraph{基于排名的扩展。}
在同样的嵌套与固定分数条件下，令 $\operatorname{rank}_{\rm P}(x)$ 为参考答案在共同主码 key 下从 1 开始的排名，则
\begin{equation}
 D_S(\alpha)=\sum_{q\in\mathcal Q_S}\sum_{x\in\mathcal R_{\rm ref}(q)}
 \mathbf1\{\operatorname{rank}_{\rm P}(x)>c(\alpha)\}.
 \label{eq:6-rank-curve}
\end{equation}
导出这些排名，可从一次参考搜索恢复整条网格曲线。当前 harness 仍逐预算搜索；排名提取与传输是同一排序条件下的后续工作。
